% !TeX spellcheck = de_DE
%\input{header.tex}

\documentclass[10pt,a4paper]{article}
\usepackage[utf8]{inputenc}
\usepackage[ngerman]{babel}
\usepackage{amsmath}
\usepackage{amsfonts}
\usepackage{amssymb}
%\usepackage{marvosym}
\usepackage{ gensymb }
%\usepackage{graphicx} % Allows for eps images
\usepackage{hyperref} %Links
\usepackage{url}
\renewcommand{\UrlFont}{\small\tt} %smaller font for urls
\usepackage[a4paper,lmargin={3cm},rmargin={3cm},tmargin={3cm},bmargin = {3cm}]{geometry}


% Literaturangabe
%\usepackage{biblatex}
%\bibliographystyle{plain}
\bibliographystyle{alpha} %http://sites.stat.psu.edu/~surajit/present/bib.htm
%\usepackage{natbib}
%\bibliographystyle{plainnat}

\usepackage{verbatim} %Fuer mehrzeilige kommentare \begin{comment} \end{comment}

%\begin{comment}
\usepackage{titlesec}
\titleformat{\subsection}[runin]% runin puts it in the same paragraph
       {\normalfont\bfseries}% formatting commands to apply to the whole heading
       {\thesubsection}% the label and number
       {0.5em}% space between label/number and subsection title
       {}% formatting commands applied just to subsection title
       []% e.g. [.] punctuation or other commands following subsection title\frac{•}{•}
%\end{comment}

%\usepackage{pdfpages} %For including PDFs
\usepackage{tikz} %For including Graphics generated by QtiPlot
%\usepackage{pdfcomment} %For generating comments in PDFs, e.g. \pdfsquarecomment{Kommentar}

%Some other shit
\usepackage{float}
\floatstyle{boxed}
\restylefloat{figure}
\usepackage{wrapfig}

\linespread{1.1}

\addto\captionsgerman{ \renewcommand{\figurename}{\small{\textbf{Abb.}}}
\captionsetup{figurewithin = section}
\captionsetup{font=small, labelfont=bf} }



\begin{document}

\pagenumbering{Roman}

%Metadaten
\title{Fortgeschrittenenpraktikum\\ Kurzversuch 6:\\ Fourier-Transform-Spektroskopie }
%\date{02.02.2015}
\author{Nico Kalogerakis \\ Léonie Lanz}

\maketitle
%\vfill
\newpage

\tableofcontents
\vfill
\newpage

\pagenumbering{arabic}

\section{Einleitung}
%Aufgabenstellung, Ziel der Messungen
Der Kurzversuch 6 - Fourier-Transform-Spektroskopie dient der Einführung in das universelle Konzept der Fourier-Transform-Spektroskopie und dem Erlernen des Programmieren einer Auswertungsroutine. Zu didaktischen Zwecken wird hier sichtbares Licht verwendet anstelle von dem gängigen Infrarotlicht.
Ziel des Experiments ist sowohl das Spektrum einer Weißlichtlampe, als auch die Absorptionsspektren von einem Interferenz-Filter und molekularem Iod zu bestimmen. Dafür wird zuerst die zeitliche Kohärenzfunktion für die Weißlichtquelle ermittelt aus der dann durch eine Fourier-Transformation auf das Spektrum der Quelle geschlossen wird. 

\section{Theoretische Grundlagen}
%Kurze Zusammenfassung der wichtigsten Überlegungen (Formeln), die zum Verständnis des Versuches notwendig sind
%Für Details: Hinweise auf die Literatur
\subsection{Weißlichtspektrum}

\subsection{Interferometrie}
Interferometrie ist das Prinzip mithilfe von Interferenzphänomenen Messungen zu vollziehen. Im Allgemeinen bezieht sich Interferometrie auf Wellen jeglicher Art, hier beschäftigen wir uns aber nur mit den elektromagnetischen - dem Licht.\cite{OptikLip}




\subsection{Kohärenz und Autokorrelation}

\subsection{Fourier-Spektroskopie}


\subsection{Formel 1} Tolle Formel $ h \rightarrow Z^0 Z^{*0} \rightarrow e^+ e^- \mu^+ \mu^- $ total inline.
\subsection{Formel 2} Wichtige Formel:
$$ a^2 = b^2 + c^2 $$
und Nummerierte Formel:
\begin{equation}
sin^2(x)+cos^2(x) = 1
\end{equation}


%\subsection{Grafik}
%\begin{figure}[h]
%\centering
%\input{img/Graph2.tex}

%\caption[Parabel]{Tolle Grafik! Sowas wie expotentiell?}
%\end{figure}
Mehr zu dem Unterschied zwischen \textit{input} und \textit{include} gibt es \href{http://de.wikibooks.org/wiki/LaTeX-W%C3%B6rterbuch:_include}{hier}





\section{Experimenteller Aufbau}
Das Fourier-Transform-Spektrometer, das in diesem Versuch verwendet wird, besteht hauptsächlich aus einem Michelson-Interferometer. Wie in \ref{•} zu sehen, werden in das Michelson-Interferometer zwei Strahlengänge gleichzeitig eingekoppelt: zum einen der zu untersuchende Weißlichtstrahlengang, zum anderen der Laserstrahlengang, der der Kalibrierung der optischen Weglängendifferenz des Interferometers dient.\\
Die Weißlichtquelle hier ist eine Weißlicht-LED und der Laser ist ein frequenzverdoppelter diode-pumped solid-state-Laser mit einer Wellenlänge von $\lambda_L=532$nm.\\\\
Das Michelson-Interferometer teilt die einfallenden Lichtbündel in zwei identische Lichtstrahlen auf, die beide über Spiegel wieder zurückgeworfen werden. Einer der beiden Spiegel des Interferometers lässt sich über einen Linearverschiebetisch angetrieben von einem Schrittmotor bewegen.\\
Nach dem  Durchlaufen des Interferometers wird der Strahlengang des Lasers auf eine einzelne Fotodiode gelenkt, während der Weißlichtstrahl erneut durch einen halbdurchlässigen Spiegel geteilt wird. Einer der beiden Weißlichtstrahlen wird direkt auf eine Fotodiode gelenkt, der andere durchläuft die zu untersuchenden Proben, bevor er auch auf eine Fotodiode trifft.\\
Alle drei messenden Photodioden sind über eine Signalverstärkung mit einem Computer verbunden.\\\\
Die Steuerung des Motors und die Aufnahme der Daten erfolgt über ein Python GUI.\\
Die kleinstmögliche Schrittgröße des Schrittmotors beträgt rechnerisch $50,5$ nm. Allerdings ist die Bewegung des Verschiebetisches nicht exakt linear, was in der Auswertung berücksichtigt werden muss. Außerdem braucht der Antrieb ein gewisses Umkehrspiel, was kompensiert werden kann, indem Änderungen der Bewegungsrichtung nur in einer guten Entfernung von dem Startpunkt der Messungen aus vorgenommen werden - hier reichen etwa 1000 Schritte.\\\\
Bevor die Messungen gestartet werden, muss die Position Null des Motors auf die Mitte des Interferenzmusters gesetzt werden.

\section{Versuchsdurchführung}
Fourier-Transform-Spektroskopie basiert darauf, die Weglängendifferenz der beiden Teilarme des Interferometers kontrolliert zu variieren. Durch die Variation der Weglängendifferenz wird auch die zeitliche Verzögerung der beiden Teilstrahlen $\Delta t$ variiert.\\\\
Der Versuch lässt sich in zwei Abschnitte aufteilen. Im ersten Messdurchlauf ist ein Interferenz-Filter die Probe, im Zweiten eine Küvette mit molekularem Iod. Das Ziel ist es die Veränderung des Weißlichtspektrums unter Einfluss der beiden Proben zu untersuchen.\\\\
Nach der Justage wird der Interferenzfilter in einen der beiden Weißlichtstrahlengänge eingebaut. Der Schlitten des Spiegels wird vom Nullpunkt bis zu Motorposition -6000 gefahren und dann 1000 Schritte zurückgefahren um das mögliche Spiel des Motors beim Umkehren des Schlittens zu verhindern.\\
Nun erfolgt eine Messung mit 10000 Schritten -- von Motorposition -5000 bis 5000 -- mit folgenden Parametern in der Software: 20 Werte pro Messpunkt, 1 Schritt pro Messpunkt und einer Abklingzeit von 10 ms.\\
Für den zweiten Messdurchlauf wird der Interferenzfilter durch die Glasküvette mit dem Iod ersetzt. An die Küvette wird eine Heizspannung von 22,5 V angesetzt, sodass die Küvettentemperatur $76,7\pm 0,1$°C beträgt. In den Referenz-Weißlichtstrahlengang wird eine leere Glasküvette eingebaut.\\
Jetzt wird der Schlitten wieder zuerst auf Position -6000, dann auf -5000 gebracht. Nun folgt wieder eine Messung nach den oben genannten Parametern.

\section{Auswertung der Messungen}
%Im Protokoll nicht unbedingt alle Messwerte einzeln auff¨uhren, sondern nur die zur Berechnung der Ergebnisse wesentlichen Daten
%Weitere Messwerte im Anhang
\subsection{Form der aufgenommenen Daten}
Das Software erzeugt pro Messdurchgang jeweils eine Datei mit $21 \times 10000$ ASCII-Daten pro Messkanal\footnote{Channel 0: Probe, Channel 1: Referenz, Channel 2: Laser}. Dabei enthält die erste Spalte die Motorpositionen und die 20 anderen Spalten die eigentlichen Messwerte, die per der Einstellung der Software pro Motorposition aufgenommen wurden.\\\\
Die Auswertung der Daten erfolgt mit einem Python Skript, das mit Visual Studio Code geschrieben wurde.\\
Nachdem die auszuwertenden Daten eingelesen wurden, werden sie in zwei Arrays aufgeteilt, die einerseits die Motorpositionen und andererseits die tatsächlichen Messungen enthalten. Dann werden die Messungen gemittelt.

\subsection{Kalibrierfunktion}
Die Kalibrierfunktion wird anhand der Laserdaten bestimmt um das nicht-lineare Verhalten des Motors zu korrigieren.\\\\
Dafür werden die Laserdaten zuerst interpoliert, um die Abtastrate zu erhöhen (siehe \autoref{fig:l_interpol}). Für die Interpolation wurde sich für eine UnivariateSpline zum 4. Grad aus der Scipy Bibliothek entschieden. Es wurde um einen Faktor von 10 interpoliert. 
\begin{figure}
\centering
\input{img\Laser Daten Interpolation vs Originaldaten.png}
\caption{Ausschnitt der Laserdaten Interpolation im Vergleich zu den Originaldaten}
\label{fig:l_interpol}
\end{figure}
Die Kalibrierfunktion ordnet jeder existierenden Motorposition $P_{ist}$ eine "Soll-Position" $P_{soll}$ zu.\\
Dafür wurden die Maxima der interpolierten Laserdaten durch die argrelextrema Funktion aus der Scipy Bibliothek bestimmt. Diese Maxima sind die $P_{ist}$, denen mithilfe einer linearen Regression ihre $P_{soll}$-Werte zugeordnet werden.\\
Die Kalibrierfunktion $\Delta(i)=P_{soll}(i)-P_{ist}(i)$ wird erst nur durch die Korrektur der Maxima berechnet und dann interpoliert, sodass sie auch auf die anderen Motorpositionen anwendbar ist.

\subsection{Berechnung des Spektrums}
Wenn die Kalibrierfunktion bestimmt ist, wird die Ortsachse der gemittelten Weißlichtdaten durch sie korrigiert. Dafür wird die Funktion auf die Positionswerte addiert.\\
Für die Fouriertransformation braucht es ein äquidistantes Gitter, welches die Probe-Messdaten und die Referenz-Messdaten gemeinsam haben müssen. Außerdem werden die gemittelten Messdaten über die korrigierte Ortsachse interpoliert und  dem erstellten Gitter zugeordnet.\\
Jetzt werden die Weglängendifferenzen $\Delta s$ in Wegzeitdifferenzen $\Delta t=\frac{\Delta s}{c}$ umgerechnet.

\section{Diskussion der physikalischen Ergebnisse}
Vergleich mit den theoretischen Erwartungen, evtl. auch mit Ergebnissen aus der Literatur

Aus den gemessenen Spektren sollen folgende Fragen beantwortet werden:
1. Wie erkl¨art sich die spektrale Verteilung der Weisslicht-LED?
2. Wie groß ist die spektrale Bandbreite des eingesetzten Absorptionsfilters und wo liegt
dessen Zentralwellenl¨ange?
3. Wie groß ist die experimentell ermittelte spektrale Aufl¨osung des Fourier-Transform-
Spektrometers? Vergleichen Sie mit der theoretischen Aufl¨osungsgrenze.

\section{Zusammenfassung}
\subsection{}  Das will ich belegen \cite[p. 2]{erstesBuch}



\newpage

%\pagenumbering{Roman}
\pagenumbering{Alph}


%\nocite{zweitesBuch}
\bibliography{praktikum.bib}
\addcontentsline{toc}{section}{Literatur}
Im Text herangezogene Quellen dort markieren und am Ende zusammenstellen


\newpage

\listoffigures %lof
\addcontentsline{toc}{section}{Abbildungsverzeichnis}
Hierhin gehören in der Regel Originalmessprotokolle (wichtig) etc. 

%\tableofcontents %toc
%\listoftables %lot

\end{document} 